\documentclass[10pt,a4paper]{amsart}
\usepackage[margin=19mm]{geometry}
\usepackage{amsmath,amssymb,mathtools}
\usepackage[scr=rsfs,cal=euler]{mathalfa}
\usepackage{xfrac}
\usepackage[unicode,hidelinks,bookmarks=false]{hyperref}
\newcommand{\ie}{i.e.\ }
\newcommand{\cf}{cf.\ }
\newcommand{\resp}{respectively\ }
\newcommand{\Subsection}{\subsection}
\providecommand{\nobreakdash}{\nobreak\hbox{-}}
\setlength{\emergencystretch}{2em}
\multlinegap0pt
\usepackage{lipsum}
\usepackage{tikz-cd}
\newtheorem{theorem}{Theorem}
\newtheorem{lemma}{Lemma}
\newtheorem{proposition}{Proposition}
\newtheorem{corollary}{Corollary}
\newtheorem{definition}{Definition}
\newtheorem{observation}{Observation}
\begin{document}
\title[The Cofree Functor on G-Sets and Its Approximations]{The Cofree Functor on G-Sets and Its Approximations}
\author{Frank Murphy-Hernandez}
\date{}
\maketitle
\noindent\emph{Source and licence.} The Cofree Functor on G-Sets and Its Approximations. Version arXiv:2608.23866v1 (CC BY 4.0). This material is distributed under the Creative Commons Attribution 4.0 International licence, http://creativecommons.org/licenses/by/4.0/, as recorded in the arXiv OAI licence metadata for this exact version and shown on the abstract page https://arxiv.org/abs/2608.23866v1. Full rights notice: licenses/arxiv-2608.23866-CC-BY-4.0.txt.\par\smallskip
\noindent\textit{Editorial scope.} Counted unit: the original \texttt{proposition} environment \texttt{prop:functoriality\_comonad} at source lines 1094--1109 together with its complete proof at lines 1111--1137; the delivered document keeps source lines 100--1138 so that every referenced label and every citation resolves inside it. Native editable source is the paper's own concatenated TeX; the AMS wrapper is editorial formatting only. No publisher class, upstream script or unrelated artwork is redistributed.
\par\smallskip\noindent\textit{Reference note.} This article ships no compiled bibliography (biblatex); the entries delivered here are composed from the author's own BibTeX (.bib) fields, with no field invented and no entry dropped.
\label{sec:cofree}

For a group $G$, we define a functor $C_G$ from the category of sets to the category of $G$-sets. This functor will be the right adjoint to the forgetful functor $U_G$.

\begin{definition}
Let $G$ be a group. For a set $X$, we define $C_G(X)$ as the set $X^G$ of all functions from $G$ to $X$. We endow $C_G(X)$ with a $G$-action given by
\[
(a\phi)(g) = \phi(a^{-1}g),
\]
for all $\phi \in X^G$ and $a, g \in G$.

For a map $\alpha \colon X \longrightarrow Y$, we define $C_G(\alpha) \colon C_G(X) \longrightarrow C_G(Y)$ by
\[
C_G(\alpha)(\phi)(g) = \alpha(\phi(g)),
\]
for all $\phi \in X^G$ and $g \in G$.
\end{definition}

\begin{proposition}\label{prop:cofree_Gset}
Let $G$ be a group and $X$ a set. Then $C_G(X)$ is a $G$-set.
\end{proposition}

\begin{proof}
Take $a, b, g \in G$ and $\phi \in X^G$. Then
\[
((ab)\phi)(g) = \phi((ab)^{-1}g) = \phi(b^{-1}a^{-1}g) = (b\phi)(a^{-1}g) = (a(b\phi))(g).
\]
Moreover,
\[
(e\phi)(g) = \phi(e^{-1}g) = \phi(eg) = \phi(g).
\]
Therefore $C_G(X)$ is a $G$-set.
\end{proof}

From the preceding proposition, it follows that $C_G$ sends sets to $G$-sets.

\begin{proposition}\label{prop:cofree_functor}
Let $G$ be a group. The assignment $C_G \colon \mathbf{Set} \longrightarrow G\text{-}\mathbf{Set}$ is functorial.
\end{proposition}

\begin{proof}
We first show that $C_G(\alpha)$ is a morphism of $G$-sets for any map $\alpha \colon X \rightarrow Y$. Let $\phi \in X^G$ and $a, g \in G$. Using the definition of the $G$-action, we have
\[
(a C_G(\alpha)(\phi))(g) = C_G(\alpha)(\phi)(a^{-1}g) = \alpha(\phi(a^{-1}g)),
\]
while
\[
C_G(\alpha)(a \phi)(g) = \alpha((a \phi)(g)) = \alpha(\phi(a^{-1}g)).
\]
Hence $(a C_G(\alpha)(\phi)) = C_G(\alpha)(a \phi)$, so $C_G(\alpha)$ is a $G$-equivariant map.

Next, we verify that $C_G$ preserves identities. For a set $X$ and $\phi \in X^G$,
\[
C_G(1_X)(\phi)(g) = 1_X(\phi(g)) = \phi(g),
\]
so $C_G(1_X)(\phi) = \phi$. Thus $C_G(1_X) = 1_{C_G(X)}$.

Finally, we check composition. Let $\alpha \colon X \rightarrow Y$ and $\beta \colon Y \rightarrow Z$ be maps. For $\phi \in X^G$ and $g \in G$,
\[
C_G(\beta\alpha)(\phi)(g) = \beta(\alpha(\phi(g))) = \beta(C_G(\alpha)(\phi)(g)) = C_G(\beta)(C_G(\alpha)(\phi))(g).
\]
Therefore $C_G(\beta\alpha) = C_G(\beta) \circ C_G(\alpha)$, which proves that $C_G$ is a functor.
\end{proof}

Having established that $C_G$ is a functor, we now prove its universal property: it is the right adjoint to the forgetful functor.

\begin{proposition}\label{prop:cofree_adjoint}
Let $G$ be a group. The functor $C_G$ is the right adjoint of the forgetful functor 
$U_G \colon G\text{-}\mathbf{Set} \longrightarrow \mathbf{Set}$, which sends each $G$-set to its underlying set.
\end{proposition}

\begin{proof}
We construct the counit $\epsilon^C \colon U_G C_G \longrightarrow 1_{\mathbf{Set}}$ as follows. 
For a set $X$, define the component 
\[
\epsilon^C_X \colon U_G(C_G(X)) \longrightarrow X
\]
by 
\[
\epsilon^C_X(\phi) = \phi(e),
\]
for $\phi \in X^G$. 
To verify naturality, let $\beta \colon X \rightarrow Y$ be a map of sets. For any $\phi \in X^G$, we have
\[
\epsilon^C_Y\bigl(U_G(C_G(\beta))(\phi)\bigr) 
= U_G(C_G(\beta))(\phi)(e) 
= \beta(\phi(e)) 
= \beta(\epsilon^C_X(\phi)).
\]
Hence the naturality square commutes, so $\epsilon^C$ is a natural transformation.

We now construct the unit $\eta^C \colon 1_{G\text{-}\mathbf{Set}} \longrightarrow C_G U_G$. 
For a $G$-set $X$, define the component 
\[
\eta^C_X \colon X \longrightarrow C_G(U_G(X)) = X^G
\]
by 
\[
\eta^C_X(x)(g) = g^{-1}x,
\]
for $x \in X$ and $g \in G$. 
We first check that $\eta^C_X$ is $G$-equivariant. For $a, g \in G$ and $x \in X$,
\[
(a \eta^C_X(x))(g) = \eta^C_X(x)(a^{-1}g) = (a^{-1}g)^{-1}x = g^{-1}ax,
\]
while
\[
\eta^C_X(ax)(g) = g^{-1}(ax) = g^{-1}ax.
\]
Thus $\eta^C_X(ax) = a \eta^C_X(x)$, so $\eta^C_X$ is a morphism of $G$-sets.

Next, we verify that $\eta^C$ is natural. Let $\alpha \colon X \rightarrow Y$ be a $G$-equivariant map. 
For $x \in X$ and $g \in G$,
\[
C_G(U_G(\alpha))(\eta^C_X(x))(g) 
= \alpha\bigl(\eta^C_X(x)(g)\bigr) 
= \alpha(g^{-1}x) 
= g^{-1}\alpha(x) 
= \eta^C_Y(\alpha(x))(g).
\]
Hence $\eta^C$ is a natural transformation.

It remains to prove the triangle identities.

First, for any $G$-set $X$ and $x \in X$,
\[
\epsilon^C_{U_G(X)}\bigl(U_G(\eta^C_X)(x)\bigr) 
= \eta^C_X(x)(e) 
= e^{-1}x 
= x.
\]
Therefore 
\[
\epsilon^C U_G \circ U_G \eta^C = 1_{U_G}.
\]

Second, for any set $X$ and any $\phi \in C_G(X) = X^G$, we compute the following chain of equalities:
\[
\begin{aligned}
C_G(\epsilon^C_X)\bigl(\eta^C_{C_G(X)}(\phi)\bigr)(g)
  &= \epsilon^C_X\bigl(\eta^C_{C_G(X)}(\phi)(g)\bigr) \\
  &= \epsilon^C_X(g^{-1}\phi) \\
  &= (g^{-1}\phi)(e) \\
  &= \phi((g^{-1})^{-1}e) \\
  &= \phi(g).
\end{aligned}
\]
Thus
\[
C_G(\epsilon^C_X)\bigl(\eta^C_{C_G(X)}(\phi)\bigr) = \phi,
\]
so
\[
C_G \epsilon^C \circ \eta^C C_G = 1_{C_G}.
\]

Since both triangle identities hold, we conclude that $C_G$ is indeed the right adjoint of the forgetful functor $U_G$.
\end{proof}

We can reformulate this adjunction in terms of Hom-functors. 
For a $G$-set $X$ and a set $Y$, we have the natural isomorphism
\[
\theta_{XY} \colon \operatorname{Hom}_{\mathbf{Set}}(U_G(X), Y) 
\longrightarrow \operatorname{Hom}_{G\text{-}\mathbf{Set}}(X, C_G(Y)),
\]
which is defined for any $\alpha \in \operatorname{Hom}_{\mathbf{Set}}(U_G(X), Y)$, 
$x \in X$, and $g \in G$ by
\[
\theta_{XY}(\alpha)(x)(g) = \alpha(g^{-1}x).
\]

\begin{proposition}\label{prop:cofree_faithfulness}
Let $G$ be a group and consider the adjunction $U_G \dashv C_G$ with unit $\eta^C \colon 1_{G\text{-}\mathbf{Set}} \longrightarrow C_G U_G$ and counit $\epsilon^C \colon U_G C_G \longrightarrow 1_{\mathbf{Set}}$. The following statements hold:
\begin{enumerate}
\item The forgetful functor $U_G \colon G\text{-}\mathbf{Set} \longrightarrow \mathbf{Set}$ is faithful.
\item The cofree functor $C_G \colon \mathbf{Set} \longrightarrow G\text{-}\mathbf{Set}$ is faithful.
\item The functor $U_G$ is full if and only if $G$ is the trivial group.
\item The functor $C_G$ is full if and only if $G$ is the trivial group.
\end{enumerate}
\end{proposition}

\begin{proof}
We analyse the components of the unit and the counit.

For a $G$-set $X$, the unit component is given by
\[
\eta^C_X \colon X \longrightarrow X^G, \qquad \eta^C_X(x)(g) = g^{-1}x.
\]
If $\eta^C_X(x_1) = \eta^C_X(x_2)$, then evaluating at $g = e$ yields $x_1 = x_2$. Therefore $\eta^C_X$ is a monomorphism in $\mathbf{Set}$ for every $X$. Since the left adjoint is faithful if and only if all components of the unit are monomorphisms, we conclude that $U_G$ is faithful.

For a set $X$, the counit component is given by
\[
\epsilon^C_X \colon X^G \longrightarrow X, \qquad \epsilon^C_X(\phi) = \phi(e).
\]
For every $x \in X$, the constant function $\phi_x \in X^G$ defined by $\phi_x(g) = x$ satisfies $\epsilon^C_X(\phi_x) = x$. Thus $\epsilon^C_X$ is an epimorphism in $\mathbf{Set}$ for every $X$. Since the right adjoint is faithful if and only if all components of the counit are epimorphisms, it follows that $C_G$ is faithful.

It remains to examine the fullness conditions. The functor $U_G$ is full if and only if every component of the unit $\eta^C_X$ is a split epimorphism. If $G$ is non-trivial, choose a set $X$ with at least two elements, regarded as a trivial $G$-set. In that case, the image of $\eta^C_X$ consists only of the constant functions from $G$ to $X$, since $\eta^C_X(x)(g) = x$ for all $g \in G$. Since there exist non-constant functions in $X^G$ when $|X| \geq 2$, the map $\eta^C_X$ is not surjective, hence not an epimorphism, and therefore cannot be a split epimorphism. Consequently, $U_G$ is not full unless $G$ is trivial. Conversely, if $G$ is the trivial group, then $X^G \cong X$ and $\eta^C_X$ is the identity, so $U_G$ is actually an isomorphism of categories and, in particular, full.

Similarly, the functor $C_G$ is full if and only if every component of the counit $\epsilon^C_X$ is a split monomorphism. If $G$ is non-trivial and $X$ has at least two elements, the map $\epsilon^C_X$ is not injective, since two different functions $G \rightarrow X$ can agree at $e$. Hence $\epsilon^C_X$ is not a monomorphism, and thus not a split monomorphism. It follows that $C_G$ is not full for non-trivial $G$. If $G$ is trivial, then $\epsilon^C_X$ is again the identity, and $C_G$ is full. Therefore both adjoint functors are full exactly in the trivial case.
\end{proof}

%%%%%%%%%%%%%%%%%%%%%%%%%%%%%%%%%%%%%%%%%%%%%%%%%%%%%%%%%%%%%%%%%%%%%%%%%%%%%
\section{Burnside's Lemma}
\label{sec:burnside}

The following result characterizes the fixed points of the cofree $G$-set $C_G(X)$, which will be essential for counting orbits via Burnside's Lemma.

\begin{proposition}\label{prop:fixed_points_global}
Let $G$ be a group and $X$ a set. Then $\phi \in C_G(X)^G$ if and only if there exists $x \in X$ such that $\phi(g) = x$ for all $g \in G$.
\end{proposition}

\begin{proof}
$(\Rightarrow)$ Suppose $\phi \in C_G(X)^G$. Set $x = \phi(e)$. For any $g \in G$, since $\phi$ is fixed by the $G$-action, we have $g^{-1}\phi = \phi$. Evaluating this equality at $e \in G$ gives
\[
(g^{-1}\phi)(e) = \phi(e).
\]
But by the definition of the action on $C_G(X)$,
\[
(g^{-1}\phi)(e) = \phi((g^{-1})^{-1}e) = \phi(g).
\]
Hence $\phi(g) = \phi(e) = x$ for all $g \in G$. Therefore $\phi$ is constant.

$(\Leftarrow)$ Conversely, suppose there is $x \in X$ such that $\phi(g) = x$ for all $g \in G$. For any $a \in G$ and any $g \in G$, we have
\[
(a \phi)(g) = \phi(a^{-1}g) = x = \phi(g).
\]
Thus $a \phi = \phi$ for every $a \in G$, so $\phi \in C_G(X)^G$.
\end{proof}

We now specialize the previous characterization to the fixed points of a single group element $g \in G$. This will be particularly useful when we apply Burnside's Lemma, as it requires counting the elements fixed by each individual group element.

\begin{proposition}\label{prop:fixed_points_element}
Let $G$ be a group, $g \in G$, and $X$ a set. Then $\phi \in C_G(X)^g$ if and only if, for each coset $Ha$ with $a \in G$ and $H = \langle g \rangle$, the restriction $\phi|_{Ha}$ is constant.
\end{proposition}

\begin{proof}
$(\Rightarrow)$ Suppose $\phi \in C_G(X)^g$, i.e. $g \phi = \phi$. 
For any $a \in G$, let $x = \phi(a)$. Since $g \phi = \phi$, we also have $g^{-1} \phi = \phi$. Evaluating at $a$ gives
\[
(g^{-1}\phi)(a) = \phi(a) = x.
\]
But by definition of the action on $C_G(X)$,
\[
(g^{-1}\phi)(a) = \phi((g^{-1})^{-1}a) = \phi(g a).
\]
Hence $\phi(g a) = \phi(a) = x$. By induction, $\phi(g^n a) = x$ for all $n \ge 0$. 
For negative powers, since $g \phi = \phi$, we have $(g \phi)(a) = \phi(a)$, which gives
\[
(g \phi)(a) = \phi(g^{-1}a) = \phi(a) = x.
\]
Thus, again by induction, $\phi(g^{-n} a) = x$ for all $n \ge 0$. Therefore $\phi$ takes the constant value $x$ on every element of the coset $Ha = \{g^n a \mid n \in \mathbb{Z}\}$. Hence $\phi|_{Ha}$ is constant.

$(\Leftarrow)$ Conversely, suppose that for each coset $Ha$ with $H = \langle g \rangle$, the restriction of $\phi$ to $Ha$ is constant. 
Let $h \in G$. Then $h$ belongs to exactly one coset, say $h \in Ha$, so there exists $n \in \mathbb{Z}$ with $h = g^n a$. Since both $g^{-1}h = g^{n-1}a$ and $h = g^n a$ lie in the same coset $Ha$, and $\phi$ is constant on this coset, we have
\[
\phi(g^{-1} h) = \phi(h).
\]
Therefore,
\[
(g \phi)(h) = \phi(g^{-1} h) = \phi(h).
\]
Since $h \in G$ was arbitrary, it follows that $g \phi = \phi$, and thus $\phi \in C_G(X)^g$.
\end{proof}

The characterization of fixed points obtained above is particularly useful in conjunction with the following classical counting result, which relates the number of orbits of a finite group action to the average number of elements fixed by each group element.

\begin{proposition}[Burnside's Lemma]\label{prop:burnside}
Let $G$ be a finite group and $X$ a finite $G$-set. Then
\[
\lvert X/G \rvert = \frac{1}{\lvert G \rvert} \sum_{g \in G} \lvert X^g \rvert,
\]
where $X^g = \{ x \in X \mid g x = x \}$ denotes the set of fixed points of $g$.
\end{proposition}
\begin{proof}
This is the classical Burnside's Lemma; see for instance \cite{rotman2012introduction}.
\end{proof}

We now apply Burnside's Lemma to the cofree $G$-set $C_G(X)$. This yields a classical number-theoretic result (Fermat's little theorem) as a direct corollary.

\begin{proposition}\label{prop:fermat_prime}
Let $n$ be a natural number and $p$ a prime number. Then $p \mid n^p + (p-1)n$.
\end{proposition}

\begin{proof}
Let $X$ be a set with $n$ elements and consider the cofree $G$-set $C_{\mathbb{Z}_p}(X)$. 
Since $p$ is prime, every non-zero element $x \in \mathbb{Z}_p$ generates the whole group. 
Thus, for any such $x$, the set of fixed points of $x$ equals the set of fixed points of the entire group:
\[
C_{\mathbb{Z}_p}(X)^x = C_{\mathbb{Z}_p}(X)^{\mathbb{Z}_p}.
\]
By our previous characterization of the fixed points of the cofree $G$-set, the elements of this set are precisely the constant maps from $\mathbb{Z}_p$ to $X$. Hence 
\[
\left\lvert C_{\mathbb{Z}_p}(X)^x \right\rvert = n
\]
for every $x \neq 0$.

On the other hand, for the identity element $0 \in \mathbb{Z}_p$, we have
\[
C_{\mathbb{Z}_p}(X)^0 = C_{\mathbb{Z}_p}(X) = X^{\mathbb{Z}_p},
\]
so its cardinality is
\[
\left\lvert C_{\mathbb{Z}_p}(X)^0 \right\rvert = n^p.
\]

Applying Burnside's Lemma to the action of $\mathbb{Z}_p$ on $C_{\mathbb{Z}_p}(X)$, we obtain
\[
\left\lvert C_{\mathbb{Z}_p}(X)/\mathbb{Z}_p \right\rvert
= \frac{1}{p} \left( n^p + (p-1)n \right).
\]
Since the left-hand side is a non-negative integer, the right-hand side must also be an integer. Therefore
\[
p \mid n^p + (p-1)n,
\]
as required.
\end{proof}

The previous result for a prime $p$ generalizes naturally to prime powers $p^m$. The following proposition provides the corresponding congruence modulo $p^m$, obtained by applying Burnside's Lemma to the cyclic group $\mathbb{Z}_{p^m}$.

\begin{proposition}\label{prop:fermat_prime_power}
Let $n$ and $m$ be natural numbers, and let $p$ be a prime number. Then
\[
p^m \mid n^{p^m} + \sum_{s=1}^m (p^s - p^{s-1}) n^{p^{m-s}}.
\]
\end{proposition}

\begin{proof}
Let $X$ be a set with $n$ elements, and consider the cofree $G$-set $C_{\mathbb{Z}_{p^m}}(X)$.

Recall that in the cyclic group $\mathbb{Z}_{p^m}$, for each $s = 0, 1, \dots, m$, there is a unique subgroup of order $p^s$. Consequently, the number of elements of order exactly $p^s$ is
\[
p^s - p^{s-1}
\]
for $s = 1, \dots, m$ (the identity element has order $p^0 = 1$).

Now take an element $x \in \mathbb{Z}_{p^m}$ of order $p^s$, where $1 \le s \le m$. The subgroup generated by $x$ is $\langle x \rangle$, whose order is $p^s$. Hence its index in $\mathbb{Z}_{p^m}$ is
\[
[\mathbb{Z}_{p^m} : \langle x \rangle] = \frac{p^m}{p^s} = p^{m-s}.
\]
By the characterization of fixed points in the cofree $G$-set, $\phi \in C_{\mathbb{Z}_{p^m}}(X)^x$ if and only if $\phi$ is constant on each coset of $\langle x \rangle$. Since there are exactly $p^{m-s}$ such cosets, the number of such maps is
\[
\left\lvert C_{\mathbb{Z}_{p^m}}(X)^x \right\rvert = n^{p^{m-s}}.
\]

For the identity element $0 \in \mathbb{Z}_{p^m}$, we have $C_{\mathbb{Z}_{p^m}}(X)^0 = C_{\mathbb{Z}_{p^m}}(X) = X^{\mathbb{Z}_{p^m}}$, so
\[
\left\lvert C_{\mathbb{Z}_{p^m}}(X)^0 \right\rvert = n^{p^m}.
\]

Applying Burnside's Lemma to the action of $\mathbb{Z}_{p^m}$ on $C_{\mathbb{Z}_{p^m}}(X)$, we obtain
\[
\left\lvert C_{\mathbb{Z}_{p^m}}(X)/\mathbb{Z}_{p^m} \right\rvert
= \frac{1}{p^m} \left( n^{p^m} + \sum_{s=1}^m (p^s - p^{s-1}) n^{p^{m-s}} \right).
\]
Since the left-hand side is a non-negative integer, the right-hand side must also be an integer. Therefore
\[
p^m \mid n^{p^m} + \sum_{s=1}^m (p^s - p^{s-1}) n^{p^{m-s}},
\]
as required.
\end{proof}

We now extend the previous prime-power case to an arbitrary cyclic group of order $n$. The following classical congruence, which generalizes Fermat's little theorem, emerges naturally from the same cofree construction.

\begin{proposition}\label{prop:fermat_general}
Let $n$ and $k$ be natural numbers. Then
\[
n \mid \sum_{d \mid n} \varphi(d)\, k^{\,n/d},
\]
where $\varphi$ denotes Euler's totient function.
\end{proposition}

\begin{proof}
Let $G = \mathbb{Z}_n$ be the cyclic group of order $n$, and let $X$ be a set with $|X| = k$. We consider the cofree $G$-set $C_G(X) = X^G$.

By the characterization of fixed points obtained in Proposition~\ref{prop:fixed_points_element}, for an element $g \in G$, the set of fixed points $C_G(X)^g$ consists precisely of those functions that are constant on each coset of the subgroup $\langle g \rangle$. Since there are $[G : \langle g \rangle]$ such cosets, and each coset can be assigned independently an element of $X$, we have
\[
\left\lvert C_G(X)^g \right\rvert = |X|^{[G : \langle g \rangle]} = k^{\, n / \operatorname{ord}(g)}.
\]

Now, in the cyclic group $\mathbb{Z}_n$, for each divisor $d$ of $n$, there exist exactly $\varphi(d)$ elements of order $d$. Grouping the elements of $G$ by their order and applying Burnside's Lemma (Proposition~\ref{prop:burnside}) to the $G$-set $C_G(X)$, we obtain
\[
\left\lvert C_G(X) / G \right\rvert
= \frac{1}{n} \sum_{g \in G} \left\lvert C_G(X)^g \right\rvert
= \frac{1}{n} \sum_{d \mid n} \varphi(d)\, k^{\, n/d}.
\]

Since the left-hand side represents the number of orbits of the $G$-action on $C_G(X)$, it is necessarily an integer. Hence,
\[
n \mid \sum_{d \mid n} \varphi(d)\, k^{\, n/d}.
\]
Thus the result follows.
\end{proof}

We now consider the elementary abelian $p$-group $(\mathbb{Z}_p)^r$ instead of the cyclic $p$-group. This choice yields a different classical congruence, which is a slight variant of the previous results.

\begin{proposition}\label{prop:fermat_elementary_abelian}
Let $p$ be a prime number, $r$ a natural number, and $k$ a natural number. Then
\[
p^r \mid k^{\,p^r} + (p^r - 1)\, k^{\,p^{r-1}}.
\]
\end{proposition}

\begin{proof}
Let $G = (\mathbb{Z}_p)^r$ be the elementary abelian $p$-group of order $p^r$, and let $X$ be a set with $|X| = k$. We consider the cofree $G$-set $C_G(X) = X^G$.

We apply Burnside's Lemma (Proposition~\ref{prop:burnside}) to this $G$-set. To do so, we compute the number of fixed points $|C_G(X)^g|$ for each $g \in G$.

By the characterization of fixed points obtained in Proposition~\ref{prop:fixed_points_element}, for any $g \in G$, we have
\[
\left\lvert C_G(X)^g \right\rvert = |X|^{[G : \langle g \rangle]} = k^{\, p^r / \operatorname{ord}(g)}.
\]

Now we classify the elements of $G$ by their order. For the identity element $g=e$, we have $\operatorname{ord}(g)=1$, and therefore $\left\lvert C_G(X)^e \right\rvert = |C_G(X)| = k^{\,p^r}$. For any non-identity element $g \neq e$, since $G$ is an elementary abelian $p$-group, it follows that $\operatorname{ord}(g)=p$. Hence, $\left\lvert C_G(X)^g \right\rvert = k^{\, p^r / p} = k^{\,p^{r-1}}$, and the number of such elements in $G$ is exactly $p^r - 1$.

Summing the contributions of all elements of $G$ and applying Burnside's Lemma, we obtain that the number of orbits of the $G$-action on $C_G(X)$ is
\[
\left\lvert C_G(X) / G \right\rvert
= \frac{1}{p^r}
\left( k^{\,p^r} + (p^r - 1)\, k^{\,p^{r-1}} \right).
\]

Since the number of orbits is necessarily an integer, the numerator must be divisible by $p^r$. Therefore,
\[
p^r \mid k^{\,p^r} + (p^r - 1)\, k^{\,p^{r-1}}.
\]
This completes the proof.
\end{proof}

We now turn to a non-abelian example. Applying the cofree construction to the dihedral group yields another classical congruence, which further illustrates the versatility of Burnside's Lemma in this context.

Consider the dihedral group $D_n$ of order $2n$, defined for $n \geq 3$ as the group of symmetries of a regular $n$-gon. Algebraically, it admits the presentation
\[
D_n = \langle r, s \mid r^n = 1,\; s^2 = 1,\; srs = r^{-1} \rangle,
\]
where $r$ denotes a rotation of angle $2\pi/n$ and $s$ denotes a reflection. Thus every element of $D_n$ can be uniquely written as $r^i$ or $sr^i$, with $0 \leq i < n$. The rotations $r^i$ form a cyclic subgroup of order $n$, while the remaining $n$ elements are reflections, each of order $2$. Hence $D_n$ is non-abelian for $n \geq 3$, since $sr \neq rs$ in general. This structure makes it an excellent test case for the cofree construction: the different orders of its elements lead to distinct contributions in the fixed-point sum $\sum_{g \in D_n} |C_{D_n}(X)^g|$, which will ultimately yield a classical congruence depending on the parity of $n$.


\begin{proposition}\label{prop:dihedral_congruence}
Let $n \ge 3$ be a natural number and $k$ a natural number. Let $D_n$ be the dihedral group of order $2n$. Then
\[
2n \mid \sum_{d \mid n} \varphi(d)\, k^{\,2n/d} + n k^{\,n},
\]
where $\varphi$ denotes Euler's totient function.
\end{proposition}

\begin{proof}
Let $G = D_n$, the dihedral group of order $2n$, which consists of $n$ rotations and $n$ reflections. Let $X$ be a set with $|X| = k$, and consider the cofree $G$-set $C_G(X) = X^G$.

We apply Burnside's Lemma (Proposition~\ref{prop:burnside}) to this $G$-set. To do so, we compute the number of fixed points $|C_G(X)^g|$ for each $g \in G$.

By the characterization of fixed points obtained in Proposition~\ref{prop:fixed_points_element}, for any element $g \in G$, we have
\[
\left\lvert C_G(X)^g \right\rvert = |X|^{[G : \langle g \rangle]} = k^{\, 2n / \operatorname{ord}(g)}.
\]

We now classify the elements of $D_n$ by their order.

\begin{itemize}
\item \textbf{Rotations:} There are $n$ rotations. For each divisor $d$ of $n$, the cyclic subgroup of rotations contains exactly $\varphi(d)$ elements of order $d$. For such an element $g$, we have $\operatorname{ord}(g) = d$, hence
\[
\left\lvert C_G(X)^g \right\rvert = k^{\, 2n / d}.
\]

\item \textbf{Reflections:} There are $n$ reflections. Every reflection has order $2$. Therefore, for each reflection $g$, we have $\operatorname{ord}(g) = 2$, and consequently
\[
\left\lvert C_G(X)^g \right\rvert = k^{\, 2n / 2} = k^{\, n}.
\]
Since there are $n$ reflections, their total contribution to the sum of fixed points is $n k^{\, n}$.
\end{itemize}

Summing the contributions from all elements of $G$ and applying Burnside's Lemma, we obtain that the number of orbits of the $G$-action on $C_G(X)$ is
\[
\left\lvert C_G(X) / G \right\rvert
= \frac{1}{2n}
\left( \sum_{d \mid n} \varphi(d)\, k^{\, 2n/d} + n k^{\, n} \right).
\]

Since the number of orbits is necessarily an integer, the numerator must be divisible by $2n$. Therefore,
\[
2n \mid \sum_{d \mid n} \varphi(d)\, k^{\, 2n/d} + n k^{\, n}.
\]
This completes the proof.
\end{proof}

%%%%%%%%%%%%%%%%%%%%%%%%%%%%%%%%%%%%%%%%%%%%%%%%%%%%%%%%%%%%%%%%%%%%%%%%%%%%%
\section{Approximations to the Cofree Functor}
\label{sec:approx}

We now introduce a family of approximations to the cofree functor, parameterized by subgroups of $G$. These generalized cofree $G$-sets will be useful when studying induced actions.

\begin{definition}\label{def:cofree_induced}
Let $G$ be a group and $H$ a subgroup of $G$. For a set $X$, we define $C_H^G(X)$ as the set of all functions from the coset space $G/H$ to $X$.

We endow $C_H^G(X)$ with a $G$-action defined by
\[
(g \phi)(aH) = \phi(g^{-1}aH)
\]
for all $\phi \in C_H^G(X)$, $aH \in G/H$, and $g \in G$. This action is well-defined and makes $C_H^G(X)$ a $G$-set.

For a map $f \colon X \longrightarrow Y$, we define the induced map
\[
C_H^G(f) \colon C_H^G(X) \longrightarrow C_H^G(Y)
\]
by
\[
C_H^G(f)(\phi) = f \circ \phi,
\]
for all $\phi \in C_H^G(X)$.
\end{definition}

When $H = \{e\}$ is the trivial subgroup, we have $G/H \cong G$, so the definition above coincides with the original cofree functor $C_{\{e\}}^G(X) \cong C_G(X)$. Thus the original cofree functor is a special case of this more general construction.

Having defined the generalized cofree $G$-set $C_H^G(X)$, we now verify that this construction is functorial in the set variable $X$.

\begin{proposition}\label{prop:cofree_induced_functor}
Let $G$ be a group and $H$ a subgroup of $G$. Then the assignment
\[
C_H^G \colon \mathbf{Set} \longrightarrow G\text{-}\mathbf{Set}
\]
is functorial.
\end{proposition}

\begin{proof}
First, let $f \colon X \longrightarrow Y$ be a map of sets. We verify that $C_H^G(f) \colon C_H^G(X) \longrightarrow C_H^G(Y)$ is a morphism of $G$-sets. For any $\phi \in C_H^G(X)$, $aH \in G/H$, and $g \in G$, we have
\[
\bigl(g C_H^G(f)(\phi)\bigr)(aH)
= C_H^G(f)(\phi)(g^{-1}aH)
= f\bigl(\phi(g^{-1}aH)\bigr),
\]
while
\[
C_H^G(f)(g \phi)(aH)
= f\bigl((g \phi)(aH)\bigr)
= f\bigl(\phi(g^{-1}aH)\bigr).
\]
Hence $g C_H^G(f)(\phi) = C_H^G(f)(g \phi)$, so $C_H^G(f)$ is $G$-equivariant.

Next, for the identity map $1_X \colon X \longrightarrow X$, and for any $\phi \in C_H^G(X)$, we have
\[
C_H^G(1_X)(\phi) = 1_X \circ \phi = \phi.
\]
Thus $C_H^G(1_X) = 1_{C_H^G(X)}$.

Finally, let $f \colon X \longrightarrow Y$ and $g \colon Y \longrightarrow Z$ be maps of sets. For any $\phi \in C_H^G(X)$,
\[
C_H^G(g \circ f)(\phi) = (g \circ f) \circ \phi = g \circ (f \circ \phi) = C_H^G(g)\bigl(C_H^G(f)(\phi)\bigr).
\]
Therefore $C_H^G(g \circ f) = C_H^G(g) \circ C_H^G(f)$.

Since all three conditions are satisfied, $C_H^G$ is a functor.
\end{proof}

The functors $C_G$ and $C_H^G$ are closely related, but $C_H^G$ cannot be a right adjoint to the full forgetful functor $U_G \colon G\text{-}\mathbf{Set} \rightarrow \mathbf{Set}$. 
Besides the uniqueness of adjoints (since $C_G$ is already a right adjoint to $U_G$), a direct attempt to construct the natural isomorphism reveals a fundamental well-definedness issue.

Indeed, suppose we try to build a natural isomorphism
\[
\theta^H_{XY} \colon \operatorname{Hom}_{\mathbf{Set}}(U_G(X), Y) 
\longrightarrow \operatorname{Hom}_{G\text{-}\mathbf{Set}}(X, C_H^G(Y))
\]
for a $G$-set $X$ and a set $Y$, mimicking the cofree case. A natural candidate would be
\[
\theta^H_{XY}(\alpha)(x)(gH) = \alpha(g^{-1}x),
\]
for $\alpha \in \operatorname{Hom}_{\mathbf{Set}}(U_G(X), Y)$, $x \in X$, and $gH \in G/H$. However, this map is not well-defined in general. If $gH = g'hH$ for some $h \in H$ (i.e., $g' = gh$), we would need
\[
\alpha(g^{-1}x) = \alpha((gh)^{-1}x) = \alpha(h^{-1}g^{-1}x).
\]
For this equality to hold for every arbitrary map $\alpha \colon X \rightarrow Y$, we must have $g^{-1}x = h^{-1}g^{-1}x$ for all $x \in X$, which is equivalent to $h x = x$ for all $h \in H$. In other words, $H$ must act trivially on $X$. Since this is not true for arbitrary $G$-sets, the adjunction fails.

Motivated by this obstruction, we introduce a restricted class of $G$-sets that arise naturally from the transitive $G$-set $G/H$.

\begin{definition}\label{def:H_generated}
Let $G$ be a group and $H$ a subgroup of $G$. We say that a $G$-set $X$ is \emph{$H$-generated} if there exists a $G$-equivariant epimorphism
\[
\coprod_{i \in I} G/H \longrightarrow X
\]
for some indexing set $I$. Equivalently, $X$ is a quotient of a coproduct of copies of the transitive $G$-set $G/H$. We denote the full subcategory of $H$-generated $G$-sets by $G\text{-}\mathbf{Set}_H$.
\end{definition}

 

We now characterize the $H$-generated $G$-sets among those that are disjoint unions of transitive $G$-sets. The condition is naturally expressed in terms of conjugacy of subgroups.

\begin{proposition}\label{prop:H_generated_criterion}
Let $G$ be a group and let
\[
X = \bigsqcup_{i \in I} G/H_i
\]
be a disjoint union of transitive $G$-sets. Then $X$ is $H$-generated if and only if for each $i \in I$, the subgroup $H$ is subconjugate to $H_i$ in $G$, that is, there exists an element $g_i \in G$ such that
\[
g_i^{-1} H g_i \subseteq H_i.
\]
\end{proposition}

\begin{proof}
$(\Rightarrow)$ Suppose $X$ is $H$-generated. Then there exists a $G$-equivariant epimorphism
\[
\pi \colon \coprod_{j \in J} G/H \longrightarrow X,
\]
for some indexing set $J$. Fix an index $i \in I$. Since $G/H_i$ is a connected component of $X$, and $\pi$ is surjective, there must be some index $j \in J$ such that the image of the corresponding copy of $G/H$ under $\pi$ is exactly $G/H_i$ (because the image of a transitive $G$-set under a $G$-equivariant map is a single orbit). Thus we obtain a $G$-equivariant epimorphism
\[
\pi_{ij} \colon G/H \longrightarrow G/H_i.
\]

Now, any $G$-equivariant map $f \colon G/H \rightarrow G/H_i$ is determined by the image of the coset $H$. Let $f(H) = g H_i$ for some $g \in G$. By equivariance, for any $h \in H$,
\[
g H_i = f(H) = f(hH) = h f(H) = h g H_i.
\]
Hence $h g H_i = g H_i$, which implies $g^{-1} h g \in H_i$ for all $h \in H$. Therefore $g^{-1} H g \subseteq H_i$, so $H$ is subconjugate to $H_i$. Since $i$ was arbitrary, the condition holds for all $i$.

$(\Leftarrow)$ Conversely, suppose that for each $i \in I$ there exists $g_i \in G$ with $g_i^{-1} H g_i \subseteq H_i$. Then the map
\[
\pi_i \colon G/H \longrightarrow G/H_i, \qquad aH \longmapsto a g_i H_i
\]
is well-defined and $G$-equivariant (because $g_i^{-1} H g_i \subseteq H_i$). Moreover, it is surjective since $G$ acts transitively on $G/H_i$.

Taking the coproduct of all these maps over $i \in I$, we obtain a $G$-equivariant map
\[
\pi \colon \coprod_{i \in I} G/H \longrightarrow \bigsqcup_{i \in I} G/H_i = X,
\]
which is surjective because each component map is. Hence $X$ is $H$-generated.
\end{proof}

The obstruction to the adjunction of $C_H^G$ motivates the study of a related concept: the $H$-trace of a $G$-set. This will allow us to isolate the part of a $G$-set that is "visible" from the transitive $G$-set $G/H$.

\begin{definition}\label{def:H_trace}
Let $G$ be a group and $H$ a subgroup of $G$. For a $G$-set $X$, we define its \emph{$H$-trace} as
\[
\operatorname{Tr}_H(X) = \bigcup \bigl\{ \operatorname{im}(\alpha) \mid \alpha \in \operatorname{Hom}_{G\text{-}\mathbf{Set}}(G/H, X) \bigr\}.
\]
In other words, $\operatorname{Tr}_H(X)$ is the union of the images of all $G$-equivariant maps from the transitive $G$-set $G/H$ into $X$.
\end{definition}

\begin{proposition}\label{prop:tr_preradical}
Let $G$ be a group and $H$ a subgroup of $G$. Then $\operatorname{Tr}_H$ is a preradical on the category $G\text{-}\mathbf{Set}$.
\end{proposition}

\begin{proof}
We must verify the two defining properties of a preradical.

First, we show that $\operatorname{Tr}_H(X)$ is a $G$-subobject of $X$. Let $X$ be a $G$-set, take $x \in \operatorname{Tr}_H(X)$, and let $g \in G$. By definition of the trace, there exists a $G$-equivariant map $\alpha \colon G/H \longrightarrow X$ and an element $aH \in G/H$ such that $\alpha(aH) = x$. Since $\alpha$ is $G$-equivariant,
\[
g x = g \alpha(aH) = \alpha(g aH).
\]
Thus $g x \in \operatorname{im}(\alpha) \subseteq \operatorname{Tr}_H(X)$. Therefore $\operatorname{Tr}_H(X)$ is a $G$-subset of $X$.

Second, we verify functoriality under morphisms. Let $\beta \colon X \longrightarrow Y$ be a $G$-equivariant map, and let $x \in \operatorname{Tr}_H(X)$. Again, there exists a $G$-equivariant map $\alpha \colon G/H \longrightarrow X$ such that $x \in \operatorname{im}(\alpha)$. Then the composition
\[
\beta \circ \alpha \colon G/H \longrightarrow Y
\]
is also $G$-equivariant, and
\[
\beta(x) = \beta(\alpha(aH)) = (\beta \circ \alpha)(aH) \in \operatorname{im}(\beta \circ \alpha) \subseteq \operatorname{Tr}_H(Y).
\]
Hence $\beta(\operatorname{Tr}_H(X)) \subseteq \operatorname{Tr}_H(Y)$. This proves that $\operatorname{Tr}_H$ is a preradical on $G\text{-}\mathbf{Set}$.
\end{proof}

Observe that $\operatorname{Tr}_H(X)$ is precisely the largest $H$-generated $G$-subobject of $X$. Indeed, if $X$ is itself $H$-generated, then $\operatorname{Tr}_H(X) = X$.

Having established that $\operatorname{Tr}_H$ is a preradical, we now show that it is idempotent. This means that applying the trace operation twice yields no new elements; the $H$-trace of an $H$-trace is already the full $H$-trace.

\begin{proposition}\label{prop:tr_idempotent}
Let $G$ be a group and $H$ a subgroup of $G$. Then the preradical $\operatorname{Tr}_H$ is idempotent. That is, for every $G$-set $X$,
\[
\operatorname{Tr}_H\bigl(\operatorname{Tr}_H(X)\bigr) = \operatorname{Tr}_H(X).
\]
\end{proposition}

\begin{proof}
Let $X$ be a $G$-set and set $Y = \operatorname{Tr}_H(X)$. We must show that $\operatorname{Tr}_H(Y) = Y$.

First, we prove the inclusion $Y \subseteq \operatorname{Tr}_H(Y)$. Take $x \in Y$. By the definition of the $H$-trace, there exists a $G$-equivariant map $\alpha \colon G/H \longrightarrow X$ and an element $aH \in G/H$ such that $\alpha(aH) = x$. Since $\operatorname{im}(\alpha) \subseteq \operatorname{Tr}_H(X) = Y$, we may consider the corestriction
\[
\overline{\alpha} \colon G/H \longrightarrow Y, \qquad \overline{\alpha}(gH) = \alpha(gH).
\]
This map is well-defined and remains $G$-equivariant because it is the same formula as $\alpha$ with its codomain restricted to $Y$. Moreover,
\[
\overline{\alpha}(aH) = \alpha(aH) = x.
\]
Thus $x \in \operatorname{im}(\overline{\alpha}) \subseteq \operatorname{Tr}_H(Y)$. Hence $Y \subseteq \operatorname{Tr}_H(Y)$.

Conversely, since $\operatorname{Tr}_H$ is a preradical, we already know from Proposition~\ref{prop:tr_preradical} that $\operatorname{Tr}_H(Y)$ is a $G$-subset of $Y$. Therefore $\operatorname{Tr}_H(Y) \subseteq Y$.

Combining both inclusions, we obtain
\[
\operatorname{Tr}_H\bigl(\operatorname{Tr}_H(X)\bigr) = \operatorname{Tr}_H(Y) = Y = \operatorname{Tr}_H(X).
\]
Thus $\operatorname{Tr}_H$ is idempotent.
\end{proof}

The following result provides a fundamental criterion: a $G$-set is $H$-generated precisely when its $H$-trace is the whole set.

\begin{proposition}\label{prop:H_generated_trace_criterion}
Let $G$ be a group, $H$ a subgroup of $G$, and $X$ a $G$-set. Then $X$ is $H$-generated if and only if $\operatorname{Tr}_H(X) = X$.
\end{proposition}

\begin{proof}
$(\Rightarrow)$ Suppose $X$ is $H$-generated. Then there exists a set $I$ and a $G$-equivariant epimorphism
\[
\pi \colon \coprod_{i \in I} G/H \longrightarrow X.
\]
Let $x \in X$ be arbitrary. Since $\pi$ is surjective, there exist $i \in I$ and $gH \in G/H$ such that
\[
\pi(\iota_i(gH)) = x,
\]
where $\iota_i \colon G/H \longrightarrow \coprod_{i \in I} G/H$ is the $i$-th canonical inclusion of the coproduct.
Now, the composition $\pi \circ \iota_i \colon G/H \longrightarrow X$ is a $G$-equivariant map, and its image contains $x$. Hence, by the definition of the $H$-trace,
\[
x \in \operatorname{im}(\pi \circ \iota_i) \subseteq \operatorname{Tr}_H(X).
\]
Since $x \in X$ was arbitrary, we have $X \subseteq \operatorname{Tr}_H(X)$. The reverse inclusion $\operatorname{Tr}_H(X) \subseteq X$ holds trivially by the definition of the trace. Therefore $\operatorname{Tr}_H(X) = X$.

$(\Leftarrow)$ Conversely, suppose $\operatorname{Tr}_H(X) = X$. For each $x \in X$, choose a $G$-equivariant map
\[
\alpha_x \colon G/H \longrightarrow X
\]
such that $x \in \operatorname{im}(\alpha_x)$; such a map exists precisely because $x \in \operatorname{Tr}_H(X)$.

By the universal property of the coproduct, the family $\{\alpha_x\}_{x \in X}$ induces a unique $G$-equivariant map
\[
\alpha \colon \coprod_{x \in X} G/H \longrightarrow X.
\]
We claim that $\alpha$ is surjective. Indeed, for any $x \in X$, there exists some $gH \in G/H$ such that $\alpha_x(gH) = x$. Therefore,
\[
\alpha(\iota_x(gH)) = \alpha_x(gH) = x,
\]
so $x$ lies in the image of $\alpha$. Hence $\alpha$ is a $G$-equivariant epimorphism. This proves that $X$ is $H$-generated.
\end{proof}

We now show that the subcategory of $H$-generated $G$-sets is a coreflective subcategory of the category of all $G$-sets, with the $H$-trace functor serving as the coreflector.

\begin{proposition}\label{prop:tr_coreflector}
Let $G$ be a group and $H$ a subgroup of $G$. Then the full subcategory $G\text{-}\mathbf{Set}_H$ of $H$-generated $G$-sets is coreflective in $G\text{-}\mathbf{Set}$. The coreflector is the $H$-trace functor $\operatorname{Tr}_H$.
\end{proposition}

\begin{proof}
We must verify the universal property of the coreflection. For every $G$-set $Y$, we need to show that the inclusion
\[
\iota_{\operatorname{Tr}_H(Y)} \colon \operatorname{Tr}_H(Y) \hookrightarrow Y
\]
has the following property: for every $H$-generated $G$-set $X$ and every $G$-equivariant map $\alpha \colon X \longrightarrow Y$, there exists a unique $G$-equivariant map
\[
\overline{\alpha} \colon X \longrightarrow \operatorname{Tr}_H(Y)
\]
such that $\alpha = \iota_{\operatorname{Tr}_H(Y)} \circ \overline{\alpha}$.

Since $X$ is $H$-generated, by Proposition~\ref{prop:H_generated_trace_criterion} we have $\operatorname{Tr}_H(X) = X$. Moreover, since $\operatorname{Tr}_H$ is a preradical (Proposition~\ref{prop:tr_preradical}), the map $\alpha$ restricts to a $G$-equivariant map
\[
\operatorname{Tr}_H(\alpha) \colon \operatorname{Tr}_H(X) \longrightarrow \operatorname{Tr}_H(Y).
\]
But $\operatorname{Tr}_H(X) = X$, so we obtain a $G$-equivariant map
\[
\overline{\alpha} := \operatorname{Tr}_H(\alpha) \colon X \longrightarrow \operatorname{Tr}_H(Y).
\]

We claim that this map satisfies the desired factorization. Indeed, for any $x \in X$,
\[
\iota_{\operatorname{Tr}_H(Y)}(\overline{\alpha}(x))
= \iota_{\operatorname{Tr}_H(Y)}(\operatorname{Tr}_H(\alpha)(x))
= \alpha(x),
\]
where the last equality follows from the definition of $\operatorname{Tr}_H(\alpha)$ as the restriction of $\alpha$ to the trace. Hence $\alpha = \iota_{\operatorname{Tr}_H(Y)} \circ \overline{\alpha}$.

It remains to prove uniqueness. Suppose there is another $G$-equivariant map
\[
\beta \colon X \longrightarrow \operatorname{Tr}_H(Y)
\]
such that $\alpha = \iota_{\operatorname{Tr}_H(Y)} \circ \beta$. Since $\iota_{\operatorname{Tr}_H(Y)}$ is a monomorphism (being an inclusion of a subset), we can cancel it on the left:
\[
\iota_{\operatorname{Tr}_H(Y)} \circ \overline{\alpha}
= \iota_{\operatorname{Tr}_H(Y)} \circ \beta
\quad \Longrightarrow \quad
\overline{\alpha} = \beta.
\]
Thus the factorization is unique.

Therefore, for every $Y$, the inclusion $\operatorname{Tr}_H(Y) \hookrightarrow Y$ is a coreflection morphism. This proves that $G\text{-}\mathbf{Set}_H$ is a coreflective subcategory of $G\text{-}\mathbf{Set}$, with coreflector $\operatorname{Tr}_H$.
\end{proof}

For normal subgroups, the set of cosets $G/H$ admits a well-defined right action of $G$, which will be useful when studying the approximations to the cofree functor and their restrictions.

\begin{proposition}\label{prop:normal_right_action}
Let $G$ be a group and $H$ a normal subgroup of $G$. If $gH = hH$, then
\[
gaH = haH \quad \text{for all } a \in G.
\]
In other words, right multiplication by any element of $G$ is well-defined on the set of left cosets $G/H$.
\end{proposition}

\begin{proof}
Assume $gH = hH$. Then $g^{-1}h \in H$. Since $H$ is normal, for any $a \in G$ we have
\[
a^{-1}(g^{-1}h)a \in H.
\]
But
\[
a^{-1}(g^{-1}h)a = (ga)^{-1}(ha).
\]
Thus $(ga)^{-1}(ha) \in H$, which implies
\[
gaH = haH.
\]
This proves the claim.
\end{proof}

For normal subgroups, the approximations $C_H^G$ are not just close to the cofree functor, but are naturally isomorphic to the $H$-trace of the cofree functor. This provides a concrete representation of the trace that will be essential for further categorical constructions.

\begin{proposition}\label{prop:CHG_iso}
Let $G$ be a group and $H$ a normal subgroup of $G$. Then
\[
C_H^G \cong \operatorname{Tr}_H \circ C_G,
\]
as functors from $\mathbf{Set}$ to $G\text{-}\mathbf{Set}$.
\end{proposition}

\begin{proof}
For each set $X$, we define a map
\[
\theta_X \colon C_H^G(X) \longrightarrow \operatorname{Tr}_H(C_G(X))
\]
by
\[
\theta_X(\phi)(g) = \phi(gH),
\]
for $\phi \in C_H^G(X)$ and $g \in G$.

We first verify that $\theta_X(\phi)$ indeed lies in $\operatorname{Tr}_H(C_G(X))$. Let $\psi = \theta_X(\phi)$. Define a map
\[
\alpha \colon G/H \longrightarrow C_G(X), \qquad \alpha(aH) = a \psi,
\]
where the action on the right is the usual $G$-action on $C_G(X)$, i.e., $(a \psi)(g) = \psi(a^{-1}g)$.

We claim that $\alpha$ is well-defined. Suppose $aH = bH$. Then $a^{-1}b \in H$. For any $g \in G$, we have
\[
(a \psi)(g) = \psi(a^{-1}g) = \phi(a^{-1}gH).
\]
Since $H$ is normal, by Proposition~\ref{prop:normal_right_action}, the equality $aH = bH$ implies that $a^{-1}gH = b^{-1}gH$ for all $g \in G$. Hence
\[
\phi(a^{-1}gH) = \phi(b^{-1}gH) = \psi(b^{-1}g) = (b \psi)(g).
\]
Thus $a \psi = b \psi$, so $\alpha$ is well-defined. Moreover, $\alpha$ is $G$-equivariant because for any $x \in G$ and $aH \in G/H$,
\[
\alpha(xaH) = (xa) \psi = x (a \psi) = x \alpha(aH).
\]
Since $\alpha(H) = \psi$, we have $\psi \in \operatorname{im}(\alpha) \subseteq \operatorname{Tr}_H(C_G(X))$. Hence $\theta_X$ is well-defined.

Next, we prove that $\theta_X$ is $G$-equivariant. For $\phi \in C_H^G(X)$, $a, g \in G$,
\[
\theta_X(g \phi)(g) = (g \phi)(gH) = \phi(g^{-1}gH) = \theta_X(\phi)(g^{-1}g) = (g \theta_X(\phi))(g).
\]
Thus $\theta_X(g \phi) = g \theta_X(\phi)$.

We now show that $\theta_X$ is injective. If $\theta_X(\phi) = \theta_X(\psi)$, then for any $gH \in G/H$,
\[
\phi(gH) = \theta_X(\phi)(g) = \theta_X(\psi)(g) = \psi(gH),
\]
so $\phi = \psi$.

It remains to prove surjectivity. Let $\psi \in \operatorname{Tr}_H(C_G(X))$. By definition, there exists a $G$-equivariant map $\alpha \colon G/H \longrightarrow C_G(X)$ such that $\psi \in \operatorname{im}(\alpha)$. Pick $aH \in G/H$ with $\psi = \alpha(aH)$. Define
\[
\phi \colon G/H \longrightarrow X, \qquad \phi(gH) = \psi(g).
\]
We must show that $\phi$ is well-defined. Suppose $gH = hH$. Then, using the equivariance of $\alpha$,
\[
g \psi = g \alpha(aH) = \alpha(gaH),
\]
and similarly
\[
h \psi = \alpha(haH).
\]
By Proposition~\ref{prop:normal_right_action}, since $gH = hH$, we have $gaH = haH$. Hence $g \psi = h \psi$. Evaluating these equal functions at the identity element $e \in G$, we obtain
\[
\psi(g) = (g \psi)(e) = (h \psi)(e) = \psi(h).
\]
Thus $\phi(gH) = \phi(hH)$, so $\phi$ is well-defined. Finally, for any $g \in G$,
\[
\theta_X(\phi)(g) = \phi(gH) = \psi(g),
\]
so $\theta_X(\phi) = \psi$. Hence $\theta_X$ is surjective, and therefore bijective.

Finally, we verify that $\theta$ is a natural transformation. Let $f \colon X \longrightarrow Y$ be a map of sets. For $\phi \in C_H^G(X)$ and $g \in G$,
\[
\theta_Y(C_H^G(f)(\phi))(g)
= C_H^G(f)(\phi)(gH)
= f(\phi(gH))
= f(\theta_X(\phi)(g))
= C_G(f)(\theta_X(\phi))(g)
= \operatorname{Tr}_H(C_G(f))(\theta_X(\phi))(g),
\]
where the last equality holds because $\operatorname{Tr}_H$ is a functor (Proposition~\ref{prop:tr_preradical}) and restricts the image of $C_G(f)$ to the trace. Thus the required naturality square commutes.

Therefore, $\theta$ is a natural isomorphism between the functors $C_H^G$ and $\operatorname{Tr}_H \circ C_G$.
\end{proof}

We now combine all the previous results to prove the main theorem of this section: the generalized cofree functor $C_H^G$ is exactly the right adjoint to the forgetful functor on the category of $H$-generated $G$-sets.

\begin{proposition}\label{prop:CHG_adjoint}
Let $G$ be a group and $H$ a normal subgroup of $G$. Then the corestricted functor
\[
C_H^G \colon \mathbf{Set} \longrightarrow G\text{-}\mathbf{Set}_H
\]
(which exists by Proposition~\ref{prop:CHG_iso}) is the right adjoint of the underlying functor
\[
U_H \colon G\text{-}\mathbf{Set}_H \longrightarrow \mathbf{Set},
\]
which forgets the $G$-action.
\end{proposition}

\begin{proof}
Let
\[
U \colon G\text{-}\mathbf{Set} \longrightarrow \mathbf{Set}
\]
denote the usual forgetful functor, and let
\[
I_H \colon G\text{-}\mathbf{Set}_H \longrightarrow G\text{-}\mathbf{Set}
\]
be the inclusion functor of the full subcategory of $H$-generated $G$-sets.

By Proposition~\ref{prop:cofree_adjoint}, we have the adjunction
\[
U \dashv C_G,
\]
where $C_G \colon \mathbf{Set} \longrightarrow G\text{-}\mathbf{Set}$ is the cofree functor.

Furthermore, by Proposition~\ref{prop:tr_coreflector}, the subcategory $G\text{-}\mathbf{Set}_H$ is coreflective in $G\text{-}\mathbf{Set}$, with coreflector $\operatorname{Tr}_H$. Hence we have the adjunction
\[
I_H \dashv \operatorname{Tr}_H,
\]
where $\operatorname{Tr}_H \colon G\text{-}\mathbf{Set} \longrightarrow G\text{-}\mathbf{Set}_H$ is the $H$-trace functor.

Composing these two adjunctions, we obtain
\[
U \circ I_H \dashv \operatorname{Tr}_H \circ C_G.
\]

Now observe that $U \circ I_H = U_H$, since both functors first forget the $G$-action and then restrict to the subcategory. Therefore,
\[
U_H \dashv \operatorname{Tr}_H \circ C_G.
\]

Finally, by Proposition~\ref{prop:CHG_iso}, we have a natural isomorphism
\[
C_H^G \cong \operatorname{Tr}_H \circ C_G.
\]
Thus we may replace the right adjoint $\operatorname{Tr}_H \circ C_G$ by the naturally isomorphic functor $C_H^G$ (corestricted to $G\text{-}\mathbf{Set}_H$). Consequently,
\[
U_H \dashv C_H^G.
\]
This proves that $C_H^G$ is the right adjoint of the underlying functor $U_H$.
\end{proof}

%%%%%%%%%%%%%%%%%%%%%%%%%%%%%%%%%%%%%%%%%%%%%%%%%%%%%%%%%%%%%%%%%%%%%%%%%%%%%
\section{Comonads and the classification of groups}
\label{sec:comonads}

In this final section we study the comonad induced on the category of sets by the cofree adjunction. Our aim is to show that the assignment $G \mapsto K_G$ is a fully faithful contravariant functor from the category of groups to the category of comonads on $\mathbf{Set}$. In particular, the cofree comonad determines the group up to isomorphism, and the essential image forms a full subcategory equivalent to $\mathbf{Grp}^{\mathrm{op}}$.

 
We begin by recalling the standard construction. Let $\mathcal C$ and $\mathcal D$ be categories, and let $F \colon \mathcal C \rightarrow \mathcal D$ and $U \colon \mathcal D \rightarrow \mathcal C$ be functors with $F \dashv U$.  Denote the unit by $\eta \colon 1_{\mathcal C} \Rightarrow U F$ and the counit by $\varepsilon \colon F U \Rightarrow 1_{\mathcal D}$. Then the endofunctor $T := U F \colon \mathcal C \rightarrow \mathcal C$ carries a canonical comonad structure, called the \emph{comonad induced by the adjunction}. Its counit is simply $\epsilon := \varepsilon \colon T \Rightarrow 1_{\mathcal C}$, and its comultiplication is given by
\[
\delta := U \eta F \colon T \Rightarrow T^2,
\]
with components $\delta_X = U(\eta_{F(X)})$. The coassociativity and counit axioms follow directly from the triangle identities of the adjunction.


We now apply this general construction to the adjunction $U_G \dashv C_G$ studied in Section~\ref{sec:cofree}. The induced comonad on $\mathbf{Set}$ will be denoted by
\[
K_G := U_G C_G.
\]
For a set $X$, we have $K_G(X) = U_G(C_G(X)) = X^G$, the set of all functions from $G$ to $X$.

\begin{definition}
For a group $G$, the comonad $K_G = (K_G, \epsilon^G, \delta^G)$ on $\mathbf{Set}$ is defined as follows. The endofunctor $K_G$ sends a set $X$ to the set of functions $X^G$, and for a map $f: X \rightarrow Y$, it acts by post-composition: $K_G(f)(\phi) = f \circ \phi$ for each $\phi \in X^G$. The counit $\epsilon_X^G: X^G \rightarrow X$ is given by evaluation at the identity element of $G$, that is, $\epsilon_X^G(\phi) = \phi(e_G)$ for all $\phi \in X^G$. Finally, the comultiplication $\delta_X^G: X^G \rightarrow (X^G)^G$ is defined pointwise by
\[
\delta_X^G(\phi)(g)(h) = \phi(gh),
\]
for all $\phi \in X^G$ and all $g, h \in G$.
\end{definition}

We briefly derive the formula for the comultiplication. The unit of the adjunction $U_G \dashv C_G$ at a $G$-set $Y$ is
\[
\eta^C_Y \colon Y \rightarrow Y^G, \qquad \eta^C_Y(y)(g) = g^{-1}y,
\]
where the dot denotes the action of $G$ on $Y$. 
Taking $Y = C_G(X) = X^G$, we have for $\phi \in X^G$ and $g,h \in G$:
\[
\eta^C_{C_G(X)}(\phi)(g)(h) = (g^{-1}\phi)(h) = \phi((g^{-1})^{-1}h) = \phi(g h).
\]
Since $\delta^G_X = U_G(\eta^C_{C_G(X)})$, the displayed formula follows.
 
The assignment $G \mapsto K_G$ is not merely a correspondence on objects; it extends to a functor in a natural way.


\begin{proposition}\label{prop:functoriality_comonad}
The rule $G \mapsto K_G$ defines a functor
\[
\Phi \colon \mathbf{Grp}^{\mathrm{op}} \longrightarrow \mathbf{Comonads}(\mathbf{Set}).
\]
More precisely, for a group morphism $f \colon G \longrightarrow H$, the induced comonad morphism
\[
\Phi(f) \colon K_H \Longrightarrow K_G
\]
has components
\[
\Phi(f)_X(\alpha) = \alpha \circ f,
\qquad 
\alpha \in K_H(X) = X^H,\; X \in \mathbf{Set}.
\]
\end{proposition}

\begin{proof}
We first check that $\Phi(f)$ is a natural transformation. For any map $u \colon X \rightarrow Y$ and any $\alpha \in X^H$,
\[
K_G(u)(\Phi(f)_X(\alpha)) = u \circ (\alpha \circ f) = (u \circ \alpha) \circ f = \Phi(f)_Y(K_H(u)(\alpha)),
\]
so the naturality square commutes.

Next, we verify that $\Phi(f)$ respects the comonad structure. 
For the counit, take $\alpha \in X^H$:
\[
\epsilon^G_X(\Phi(f)_X(\alpha)) = (\alpha \circ f)(e_G) = \alpha(f(e_G)) = \alpha(e_H) = \epsilon^H_X(\alpha),
\]
because $f$ is a group morphism. For the comultiplication, we need
\[
\delta^G \circ \Phi(f) = \bigl( K_G(\Phi(f)) \circ \Phi(f)_{K_H} \bigr) \circ \delta^H.
\]
Fix $X$, $\alpha \in X^H$, and $g_1,g_2 \in G$. 
The left-hand side evaluated at these elements gives
\[
\delta^G_X(\Phi(f)_X(\alpha))(g_1)(g_2) = \Phi(f)_X(\alpha)(g_1 g_2) = \alpha(f(g_1 g_2)).
\]
On the right-hand side we obtain
\[
\delta^H_X(\alpha)(f(g_1))(f(g_2)) = \alpha(f(g_1) f(g_2)).
\]
Since $f$ is a morphism, $f(g_1 g_2) = f(g_1)f(g_2)$, and hence the two sides coincide. Therefore $\Phi(f)$ is a morphism of comonads. The functoriality axioms (preservation of identities and composition) follow immediately from the definitions.
\end{proof}


\begin{thebibliography}{99}
\bibitem[]{rotman2012introduction} Rotman, Joseph J.. An Introduction to the Theory of Groups. 148 (2012)
\end{thebibliography}
\end{document}
